Automated presentation generation has entered practical use, enabling users to obtain complete slide decks from natural-language requests and optional reference materials. Such requests often involve delivery requirements, including slide count, formatting, source grounding, and mandatory content, while also requiring slides to be free of severe visual defects such as occlusion and cropping. We introduce SlideGate, a benchmark comprising 40 presentation-generation tasks across seven application scenarios and five input forms, with 2,191 task-specific verification items and a protocol that separately reports instruction following, content quality, and visual presentation. Applying SlideGate to 240 outputs from six publicly available commercial slide-generation systems, we obtain content judgments that agree with two independent blinded annotators at rates of 80.7\% and 85.6\%, respectively, and system rankings based on mean per-slide visual scores that correlate with human preference rankings at Spearman \(\rho=0.855\). Across these outputs, SlideGate identifies 124 decks with blocking failures, including hard instruction violations, missing required content, or page-level cropping or occlusion. Using scoring schemes based on the published criteria of six representative evaluation frameworks, the initial grouping analysis finds no significant score difference for five schemes between decks with and without detected blocking failures. Even PPTBench-style scoring, which exhibits a significant group difference, incurs a false-positive rate of 78.0\% when detecting approximately 90\% of failing decks in leave-one-task-out validation. These findings highlight the value of combining task-specific verification with explicit failure diagnosis to reveal delivery issues that holistic scores may obscure.
